\section{Preserving possible views does not yet choose a common policy}
\label{sim:temporal-family-policies}

The correction interval leaves two questions. A calculation must first preserve every information view consistent with the supplied clocks. The decision maker must then choose instructions that can actually be followed when the relevant view may itself be unknown. Keeping both policy tables answers the first question. It does not, by itself, answer the second.

\subsection{Share calculations while retaining their different supports}

A complete temporal family carries the original clock restrictions, its checked partition, every nonempty branch and the source outcome on that branch. Several branches can select the same revision; several different revisions can also imply the same financial and observation problem. The numerical calculation can then be stored once, provided that all the distinct source and timing supports remain attached to it. Removing a repeated numerical answer must not remove a possible history.

The selected binding reference therefore compares complete policy sets and their full financial-row populations. A same-meaning correction can reuse the numerical result while changing the source binding. A changed financial assumption or observation scheme requires its corresponding numerical problem. Even two empty policy sets retain their different scenario meanings. Their common cardinality, zero, is not sufficient evidence that the underlying questions are interchangeable.

The request and saved family are checked against expected identities supplied outside the transported record. The request includes the actor, original sources, clock observations, caller's admitted evidence and computation budget. The family retains the full temporal proof and each unresolved outcome. These checks detect alteration and missing support relative to that supplied record. They do not authenticate a document or establish that the actor actually received it.

A budget stop needs special care. The earlier temporal procedure does not save a partial partition that can be resumed. The binding reference therefore preserves the entire original clock domain as unresolved and supplies no completed-view policy answer. This conservative remainder keeps all possibilities, at the cost of repeating work later. Describing it as an incremental frontier would claim a computational capability that has not been implemented.

\subsection{Decisions must agree on the histories the actor can distinguish}

Let $z$ denote a possible observation mechanism and $w$ a possible financial history. At decision date $t$, the actor receives a message $m_t(z,w)$ from that mechanism. The message may differ across histories, and the meaning of a given message may differ across mechanisms. Let $\eta_t(z)$ record a mechanism identifier only when the model explicitly makes it observable by that date. Without such an observation, $\eta_t$ has the same null value for every mechanism.

For the two-date example, the information used by a deterministic policy is
\begin{align}
 h_1(z,w)&=\bigl(m_1(z,w),\eta_1(z)\bigr),\\
 h_2(z,w,a_1)&=\bigl(h_1(z,w),a_1,m_2(z,w,a_1),\eta_2(z)\bigr).
 \label{eq:common-mechanism-history}
\end{align}
The second history retains the first message and the actor's own instruction. It contains no future message, hidden financial state or internal proof identifier. The action at each date is a function of the corresponding history, $a_t=\pi_t(h_t)$. Consequently,
\begin{equation}
 h_t(z,w)=h_t(z',w')\quad\Longrightarrow\quad
 \pi_t\bigl(h_t(z,w)\bigr)=\pi_t\bigl(h_t(z',w')\bigr).
 \label{eq:cross-mechanism-nonanticipation}
\end{equation}
This is the common-policy restriction: indistinguishable available histories require the same instruction. It applies across possible mechanisms as well as across financial histories within one mechanism. Feasibility must then hold for every retained world under the selected stress question. No probability distribution is needed for that requirement.

Internal source hashes can identify $z$ for the analyst without identifying it for the actor. Putting such a hash into $h_1$ would add information and could change the answer. Conversely, if a real message identifies its source or interpretation, omitting that information would solve a different problem. The observation declaration must say which case is intended and why.

\subsection{Four policies in each view can leave no common instruction}

Consider again receipt 7, the committed facility and the refined financial problem. Its complete row population supports four policies. For a concrete test, call the first messages $L$, $H$ and $X$. Mechanism A sends these messages for, respectively, low conditions without a trigger, high conditions without a trigger, and high conditions with a trigger. Mechanism B exchanges the meanings of $H$ and $X$. These are deliberately supplied synthetic encodings; they are not an empirical claim about a company's communications.

Each known encoding preserves the original four policies by relabeling their observations. The difficulty appears when the actor sees $H$ but does not know which encoding applies. Under A, the instruction must be pay; under B, it must be draw. A common policy cannot satisfy both requirements. The two individually feasible problems therefore have no jointly feasible policy when the mechanism is hidden.

If the encoding is identified before the first decision, the two message histories are distinguishable. One of four policies can be selected for each encoding, giving $4\times4=16$ feasible joint policies. If the encoding is identified only at the second date, the earlier contradiction remains. Later learning cannot repair a first instruction that had to satisfy incompatible requirements at the time it was chosen.

The conclusion depends on when the incompatibility arises. In a separate elementary liquidity example, the only first instruction is wait. At the second date, a company with cash 1, an amount due of 1 and no facility must pay. Another possible company has cash 0, the same amount due and a facility of 1, so it must draw. Two mechanisms exchange the meanings of the messages describing those circumstances. Hidden message meaning leaves no common second instruction; revealing the mechanism at the second date restores one feasible policy because waiting had created no earlier conflict. If the second message itself has an unambiguous meaning, learning can resolve the uncertainty without revealing a mechanism's name. This example concerns stipulated cash feasibility, not the accounting recognition of a facility or a complete company report.

A different case shows why late revelation can expand a policy set even when the mechanisms have the same financial meaning. Two possible correction heads can both select the same refined problem. All four of its policies share the same first instructions; only their later choices differ. Learning the source identity at the second date then permits four continuations for each source, giving 16 joint policies. When that identity is never observed, the continuations must agree and only four policies remain. A first attempt incorrectly expected four in both cases; the complete-set comparison exposed the omitted late distinction, and the corrected independent relabeling preserves the failed attempt as evidence.

\pagebreak
\subsection{What a complete conditional answer must retain}

The original coarse/refined correction example remains different from the swapped-encoding example. The coarse member already has no feasible policy. Identifying that member does not supply the refined signal it lacks. The joint policy problem therefore stays infeasible even if the actor is told which mechanism applies. An unresolved source outcome for which no observation map is specified instead leaves the policy question unresolved; it is not another zero-policy result.

The checked finite method retains the full financial population for every mechanism and solves the shared-history constraints in two ways: enumeration of feasible instructions and an independent formulation of forbidden action combinations. It also compares the complete policy sets with explicit relabelings of the original policies. The checks distinguish a missing action row from an available safe action, and preserve the original temporal supports when equal numerical answers share storage.

These results explain how unresolved uncertainty can persist without requiring a separate explicit node for every time or transaction possibility. Symbolic temporal restrictions represent many possible times; shared numerical results avoid repeated calculation; common histories impose the information limits on decisions. Each operation preserves a different relation and needs its own justification. General observation languages, longer horizons, empirical actor access, saved partial frontiers and durable recovery require further work.\footnote{Complete-family binding, 26 family-change cases, 26 adversarial controls and saved independent replay are retained in \path{data/accounting-master/temporal-family-binding-20261027/}. The 21 cross-mechanism policy cases, independent full-set comparisons and corrected late-revelation expectation are retained in \path{data/accounting-master/cross-mechanism-policy-20261027/}. All message mappings and visibility regimes are explicit synthetic assumptions; neither reference introduces probabilities, actual actor knowledge or accounting permission.}
